\section{Troubleshooting}
Troubleshooting advice for common issues encountered during DELT-Hit analysis can be found in Table~\ref{tab:troubleshooting}.


\begin{table}[htbp]
\centering
\small
\caption{Common issues and solutions}
\label{tab:table13}
\begin{tabular}{p{2.0cm}p{3.0cm}p{3.8cm}p{5.6cm}}
\toprule
Step & Problem & Possible Cause & Solution \\
\midrule
Experimental Setup & \texttt{command not found: conda} & Conda not installed or not initialized in the shell & Install Conda as described in the Experimental Setup section and initialize it in your terminal. On macOS: \path{~/miniconda3/bin/activate}; on Windows: \path{~/miniconda3/etc/profile.d/conda.sh} \\
9 & Command fails with unrecognized arguments & Trailing spaces after \textbackslash{} in multiline commands; mistyped command name & Remove trailing spaces from multiline commands and verify the command name. Use \texttt{delt-hit \textless{}cmd\textgreater{} --help} for accepted arguments \\
9 & Demultiplexing fails: \texttt{No such file or directory} & Trailing whitespace in the file name defined in the Excel configuration sheet & Remove whitespace from the file name and regenerate the configuration with \texttt{delt-hit init} \\
13 & Low demultiplexing efficiency (\textless{}50\%) & Incorrect or inverted barcode sequences & Increase \texttt{max\_error\_rate}; validate barcode sequences against the sequencing data; check for adapter contamination \\
13 & \texttt{qc/report.txt} does not open & Text editor does not support the file encoding & Display the file in the terminal with \texttt{cat qc/report.txt} \\
Experimental Setup / 18 & R script failures & Missing R packages or R not available & Reinstall R into the conda environment; reinstall packages with \texttt{BiocManager}; verify with \texttt{which R}; R 4.1+ required \\
18 & Empty hit lists & Statistical thresholds too stringent & Adjust FDR cutoffs (0.1 or 0.2 for exploratory analysis); check replicate consistency \\
15 & Chemical structure errors & Invalid SMILES or incorrect reaction SMIRKS & Validate reaction definitions on a small subset before full enumeration; validate building block structures with RDKit; rerun \texttt{delt-hit library enumerate} with \texttt{--debug=invalid --errors=raise} to stop at the first failing combination and save its reaction graph for inspection \\
Experimental Setup & \texttt{command not found: delt-hit} & Conda environment not activated & Run \texttt{conda activate delt-hit} before executing DELT-Hit commands \\
11 & Permission denied for shell scripts & Executable permission not set & Run \texttt{chmod u+x \textless{}script\_name.sh\textgreater{}} before executing the script \\
18 & Poor replicate correlation & Batch effects or technical variability & Review the selection protocol for consistency; check for contamination \\
\bottomrule
\end{tabular}
\end{table}


\subsection{Performance optimization guidelines}
Approximate memory requirements by library size are:

\begin{itemize}
  \item Small libraries (\textless{}1M compounds): 16 GB RAM is usually sufficient
  \item Medium libraries (1--50M compounds): 32 GB RAM is recommended
  \item Large libraries (\textgreater{}50M compounds): at least 32 GB RAM is required
\end{itemize}
To improve processing speed:

\begin{itemize}
  \item use SSD storage to improve I/O performance
  \item increase the \texttt{num\_cores} parameter for parallel processing
  \item use compressed FASTQ files (\texttt{.gz}) to reduce disk usage
\end{itemize}

\subsection{Quality control thresholds}
Suggested quality indicators for demultiplexing are:

\begin{itemize}
  \item Excellent: \textgreater{}90\% initial read retention
  \item Good: 70--90\% initial retention
  \item Acceptable: 50--70\% initial retention
  \item Poor: \textless{}50\% retention (requires troubleshooting)
\end{itemize}
Suggested quality indicators for statistical analysis are:

\begin{itemize}
  \item Replicate correlation: $R^2 > 0.7$ for biological replicates and $R^2 > 0.9$ for technical replicates
  \item Positive and negative controls: controls should show the expected behavior
\end{itemize}
